\section{Telescoping Sums}

One of the most useful patterns in summation is the phenomenon of
\emph{telescoping}.  In a telescoping sum, many terms cancel with
neighboring terms, leaving only a small number of terms from the
beginning and end of the expression.

\subsection*{A Simple Example}

Consider the sum

\[
(2-1) + (3-2) + (4-3) + (5-4).
\]

Expanding the expression gives

\[
2 - 1 + 3 - 2 + 4 - 3 + 5 - 4.
\]

Most terms cancel:

\[
\cancel{2} - 1 + \cancel{3} - \cancel{2} + \cancel{4} - \cancel{3} + 5 - \cancel{4}.
\]

After cancellation, only the first and last terms remain:

\[
5 - 1.
\]

Thus

\[
(2-1) + (3-2) + (4-3) + (5-4) = 5 - 1.
\]

\subsection*{Definition}

\begin{definition}[Telescoping Sum]
A sum is called a \emph{telescoping sum} if consecutive terms cancel when
the sum is expanded, leaving only a few remaining terms.
\end{definition}

\subsection*{General Form}

A common telescoping structure is

\[
\sum_{k=a}^{b-1} (f(k+1) - f(k)).
\]

Writing out the first few terms,

\[
\begin{aligned}
& (f(a+1)-f(a)) \\
& + (f(a+2)-f(a+1)) \\
& + (f(a+3)-f(a+2)) \\
& \quad \vdots \\
& + (f(b)-f(b-1)).
\end{aligned}
\]

All intermediate terms cancel:

\[
-f(a+1) + f(a+1), \quad
-f(a+2) + f(a+2), \quad \dots
\]

leaving only

\[
f(b) - f(a).
\]

